% =====================================================================
%  DISKUSSION  (Abschnitt 1.5 der Anleitung)
%  Ergebnisse inkl. Fehler und Literaturwerten in derselben Einheit;
%  statistische vs. systematische Fehlerquellen bewerten.
% =====================================================================
\section{Diskussion}
\label{cap:Diskussion}

Die erzielten Ergebnisse werden im Vergleich mit den Literaturwerten diskutiert:
\begin{itemize}
    \item \textbf{Systematische Fehler}: reproduzierbare Abweichung in eine Richtung, z.B. falsch geeichte Messinstrumente, Kalibrierfehler. Wiederholung verbessert die Präzision nicht.
    \item \textbf{Statistische Fehler}: Schwankungen äußerer Bedingungen (Temperatur, Luftdruck), Ablesegenauigkeit, Toleranzen der Experimentatoren. Wiederholung verbessert die Genauigkeit (siehe Gleichung \eqref{equ:FehlerMittelwert}).
    \item Grundsatz: ehrliche Abschätzung; im Zweifelsfall Fehler lieber überschätzen als unterschätzen.
\end{itemize}

\noindent Der Literaturwert von ... \cite{Example2023} liegt innerhalb/außerhalb des Unsicherheitsbereichs des ermittelten Wertes ... . Dies lässt sich dadurch begründen, dass ...
